\section{Considerações e melhorias posteriores}
\label{sec:melhorias}

Este capítulo cumpre duas funções. Na primeira seção, retoma a pergunta de
pesquisa e sintetiza o que os resultados preliminares do
Capítulo~\ref{sec:resultados} permitem e não permitem afirmar. Nas seções
seguintes, organiza a agenda de melhorias posteriores em quatro frentes ---
medidas alternativas de exposição, técnicas econométricas alternativas,
tratamento do problema de medida na coleta da pesquisa e extensões de escopo ---
com a indicação, em cada caso, de qual limitação específica a melhoria pretende
resolver.

\subsection{Considerações sobre os resultados obtidos}

A pergunta de pesquisa é se trabalhadores em ocupações mais expostas à
inteligência artificial passaram a apresentar transições de trabalho distintas
após o lançamento público do ChatGPT. A resposta que os dados sustentam hoje é
qualificada em três níveis.

No nível descritivo, a resposta é que houve mudança generalizada, e não
mudança concentrada nas ocupações expostas. A permanência na mesma faixa de
exposição caiu entre 3,6 e 8,1 pontos percentuais em todas as cinco faixas, e as
margens de mobilidade ascendente e descendente subiram quase na mesma proporção.
O mercado de trabalho brasileiro ficou mais móvel no período; a exposição
ocupacional não explica esse movimento agregado.

No nível de associação condicional, a resposta é que existe um gradiente
detectável, porém modesto. Controlada a ocupação de origem, o par UF-trimestre,
a teletrabalhabilidade interagida com o período e o conjunto de covariáveis, um
desvio-padrão adicional de exposição associa-se a aproximadamente 2,2 pontos
percentuais a mais de transição para ocupação menos exposta e a 1,4 ponto
percentual a menos de transição para ocupação mais exposta. Nas margens de
vínculo o sinal é o inverso do esperado pela leitura alarmista: maior exposição
associa-se a menos saída do emprego e a menos informalização, e a informalização
do período aparece associada à teletrabalhabilidade, não à exposição.

No nível causal, a resposta é que o desenho atual não sustenta a afirmação, por
duas razões de natureza distinta. A primeira é de escopo: esta versão do
trabalho reporta um único cenário de estimação, sem estudo de evento, sem teste
de tendências pré-marco e sem bateria de sensibilidade. Um coeficiente que não
foi submetido a esses testes não pode ser apresentado como efeito causal,
qualquer que seja o seu valor-$p$. A segunda é de medida, e é a mais séria: a
taxa medida de mudança de ocupação cai pela metade durante os trimestres de
coleta telefônica de 2020 e 2021 e volta ao patamar anterior em 2021T4, uma
descontinuidade maior que qualquer coeficiente estimado aqui e que atravessa o
período pré. Enquanto ela não for tratada, todo contraste entre pré e pós que
use a série inteira é, em parte, contraste entre modos de coleta. É esse fato
que organiza toda a agenda a seguir e que promove o tratamento do problema de
coleta de item de robustez a prioridade.

\subsection{Medidas alternativas de exposição à inteligência artificial}
\label{sec:exposicoes}

A primeira frente de robustez ataca a validade da variável de tratamento. Hoje
o trabalho depende de uma única medida, o AIOE de \citeonline{felten2021},
construída para a classificação ocupacional americana, anterior à difusão dos
modelos generativos e transportada para a classificação brasileira por uma ponte
com pesos uniformes declarados. Cada uma dessas três características é uma
hipótese que pode falhar, e a maneira de testá-las é substituir a medida
mantendo tudo o mais constante.

\subsubsection{Exposição específica a modelos de linguagem}

A substituição mais direta é pela medida de \citeonline{eloundou2024},
construída tarefa a tarefa para modelos de linguagem de grande porte e
contemporânea ao marco temporal usado aqui. Ela responde à objeção mais óbvia
que o trabalho atual atrai: o AIOE mede exposição à inteligência artificial em
sentido amplo, incluindo visão computacional e reconhecimento de padrões, que
não são a tecnologia cujo lançamento define o marco. Se o gradiente estimado com
a medida específica for maior que o estimado com a medida ampla, isso é evidência
a favor da leitura pretendida; se for menor ou nulo, é evidência de que o
gradiente atual capta outra coisa.

\subsubsection{Exposição por conteúdo de tarefas rotineiras e não rotineiras}

Uma segunda substituição usa medidas de conteúdo de tarefas, na tradição de
\citeonline{acemoglu2020}, para separar exposição a automação de rotinas de
exposição a automação de tarefas cognitivas. Essa distinção é a que dá sentido
econômico à direção prevista do efeito e permite testar se o gradiente estimado
é específico à dimensão cognitiva ou se reflete simplesmente rotinização.

\subsubsection{Teletrabalhabilidade como tratamento placebo}

Dada a correlação de $0{,}702$ entre exposição e teletrabalhabilidade, um
exercício informativo é inverter os papéis: estimar a mesma especificação usando
a teletrabalhabilidade como tratamento e a exposição como controle. Se o
gradiente de teletrabalhabilidade for da mesma ordem de grandeza, a
interpretação em termos de inteligência artificial fica insustentável e o
resultado passa a ser sobre reorganização do trabalho remoto.

\subsubsection{Sensibilidade à construção da ponte ocupacional}

A ponte COD-ISCO-SOC usa pesos uniformes entre códigos SOC associados ao mesmo
código COD, comparação em dois dígitos no grupo da agropecuária e ausência de
atribuição a ocupações militares. Convém reestimar com pontes alternativas ---
pesos proporcionais ao emprego americano por SOC, comparação em três dígitos
para todos os grupos, e restrição da amostra aos códigos COD com correspondência
unívoca --- para medir quanto do resultado depende dessas escolhas de
construção.

\subsubsection{Exposição medida no destino, e não só na origem}

Por fim, todas as covariáveis e a exposição pertencem hoje à origem do par. Uma
extensão natural é modelar a exposição do destino como escolha, o que
transformaria o problema de um gradiente de transição em um modelo de escolha
ocupacional com exposição como atributo da alternativa.

\subsection{Técnicas econométricas alternativas}

A segunda frente ataca a identificação. O procedimento atual é um modelo de
probabilidade linear com interação entre tratamento contínuo e indicador de
período, estimado com efeitos fixos de alta dimensão. Essa escolha é defensável,
mas está longe de ser a única, e os problemas diagnosticados no
Capítulo~\ref{sec:resultados} pedem estimadores construídos justamente para
tratar deles.

\subsubsection{Diferenças em diferenças com tratamento contínuo}

O estimador atual impõe, implicitamente, que a resposta seja linear na dose de
exposição e que a comparação entre doses diferentes seja livre de seleção.
\citeonline{callaway2024} mostram que essa hipótese é mais forte do que parece e
propõem parâmetros e estimadores que separam o efeito médio em cada nível de dose
da comparação entre níveis. Aplicar esse arcabouço permitiria responder se o
gradiente estimado é um efeito de dose ou um artefato de composição entre
ocupações com doses diferentes.

\subsubsection{Análise formal de sensibilidade a tendências não paralelas}

O passo anterior a qualquer análise formal de sensibilidade é o estudo de
evento: substituir a interação única com o indicador de período por interações
com cada trimestre relativo ao marco, examinar se os coeficientes pré são planos
e submetê-los a teste conjunto. Feito isso, a análise de
\citeonline{rambachan2023} responde à pergunta seguinte --- qual é a menor
violação de paralelismo no período pós, expressa como múltiplo do maior desvio
pré observado, capaz de fazer o intervalo de confiança passar a incluir zero. A
melhoria consiste em executar os dois e reportar, para cada desfecho, o
intervalo robusto em vez do intervalo convencional, discutindo qual magnitude de
violação é economicamente plausível dado o que se sabe sobre a coleta da
pesquisa.

\subsubsection{Grupos de comparação construídos por pareamento}

Uma alternativa ao controle por efeitos fixos é construir explicitamente um
contrafactual: parear ocupações de alta exposição com ocupações de baixa
exposição semelhantes em trajetória pré-marco, escolaridade média, composição
setorial e teletrabalhabilidade, e estimar a diferença apenas entre pares
comparáveis. Métodos de controle sintético aplicados a séries ocupacionais
agregadas seguem a mesma lógica e têm a vantagem de tornar visível a qualidade
do ajuste pré-marco, que é exatamente o ponto fraco do desenho atual.

\subsubsection{Modelos de duração e de transição multiestado}

O horizonte de um trimestre é uma limitação de desenho, não de dados. A PNADC
permite acompanhar o mesmo indivíduo por até quatro transições consecutivas.
Modelos de duração e modelos de transição multiestado --- entre ocupado formal,
ocupado informal, desocupado e fora da força de trabalho --- usariam essa
estrutura e permitiriam separar efeitos sobre a taxa instantânea de saída de
efeitos sobre o destino da saída.

\subsubsection{Modelos de resposta binária e efeitos marginais}

Os desfechos binários são hoje estimados por modelo de probabilidade linear, por
razões de viabilidade computacional com efeitos fixos de alta dimensão. Como
robustez, convém reestimar os principais desfechos com modelo de resposta
binária em especificações com menos efeitos fixos e comparar os efeitos
marginais médios, verificando se a aproximação linear é adequada nas faixas de
probabilidade observadas.

\subsubsection{Nível de agrupamento e inferência}

O agrupamento atual é feito na unidade primária de amostragem, com mais de trinta
mil grupos. Como a exposição varia entre ocupações, e não entre UPAs, há
argumento sério para agrupar também --- ou principalmente --- na ocupação de
origem, e para reportar o mesmo coeficiente sob agrupamento em ocupação, em
ocupação e UPA e em ocupação e trimestre. Essa última estrutura reduz os graus
de liberdade a algo próximo do número de trimestres, faixa em que a aproximação
assintótica é frágil e em que procedimentos de reamostragem apropriados a poucos
grupos, com imposição da hipótese nula, devem substituir a inferência
convencional.

\subsubsection{Correção para multiplicidade}

São sete desfechos correlacionados, estimados sobre a mesma base e com a mesma
especificação, e o trabalho não aplica hoje nenhuma correção para multiplicidade
--- apenas registra a ressalva. Convém definir, de antemão, a família de
hipóteses e o critério de decisão, e reportar o resultado corrigido como
resultado principal.

\subsection{Tratamento do problema de medida na coleta}

A descontinuidade na taxa medida de mudança de ocupação entre o período de
coleta presencial e o de coleta telefônica é maior que qualquer efeito estimado
neste trabalho. Enquanto ela não for tratada, nenhuma estimativa que use o
período de 2020 e 2021 é confiável. Três encaminhamentos são possíveis, em ordem
crescente de ambição: restringir a amostra ao período de coleta homogênea, o que
custa metade dos trimestres pré; modelar explicitamente o modo de coleta como
efeito fixo interagido com a ocupação; ou tratar a mudança de modo de coleta como
um choque de medida cuja magnitude pode ser estimada e descontada.

\subsection{Extensões de escopo}

Três extensões ampliam a pergunta sem alterar o desenho.

A primeira é incorporar rendimentos. O trabalho atual mede apenas margens de
transição, e a literatura de referência --- \citeonline{brynjolfsson2025} e
\citeonline{humlum2025} --- concentra-se em produtividade e remuneração. A PNADC
publica rendimento habitual e efetivo, o que permitiria estimar o gradiente de
exposição sobre a variação de rendimento entre entrevistas consecutivas.

A segunda é a dimensão regional. Os efeitos fixos de UF interagida com trimestre
absorvem hoje toda a variação regional. Uma extensão que tratasse mercados de
trabalho locais como unidade de exposição, no espírito de
\citeonline{acemoglu2020}, transformaria essa variação absorvida em variação
identificadora.

A terceira é a margem de entrada. O desenho atual condiciona a origem a estar
ocupada, e portanto não observa efeito algum sobre quem procura o primeiro
emprego. É plausível que seja justamente aí que um efeito de exposição apareça
primeiro, e é exatamente a margem que o desenho exclui. O exame de
heterogeneidade por faixa etária entre os já ocupados, que é o teste mais
próximo disponível, também está por fazer.

\subsection{Encaminhamento}

A ordem sugerida de execução segue o princípio de atacar primeiro o que ameaça
diretamente o resultado. Primeiro, o tratamento do problema de coleta e a
definição de qual janela homogênea será adotada como cenário principal: enquanto
essa decisão não estiver tomada e justificada, nenhum coeficiente do
Capítulo~\ref{sec:resultados} é defensável como estimativa causal. Segundo, o
estudo de evento sobre a janela escolhida, com teste conjunto dos coeficientes
pré: é ele que diz se existe desenho a defender. Terceiro, a bateria mínima de
sensibilidade --- janelas alternativas, estruturas alternativas de agrupamento,
tendência específica de exposição e especificação sem o controle de
teletrabalhabilidade na mesma amostra. Quarto, a substituição da medida de
exposição pela medida específica a modelos de linguagem e o exercício de placebo
com teletrabalhabilidade: se o gradiente não sobreviver a esses dois testes, as
demais melhorias são irrelevantes. Quinto, a análise formal de sensibilidade a
tendências não paralelas e os estimadores alternativos para tratamento contínuo.
As extensões de escopo vêm por último, porque ampliam a pergunta antes de a
resposta atual estar segura.
